\documentclass[11pt,letterpaper,oneside]{report}
\usepackage[T1]{fontenc}
\usepackage[utf8]{inputenc}
\usepackage{lmodern}
\usepackage[letterpaper,margin=0.92in,headheight=15pt]{geometry}
\usepackage{amsmath,amssymb,amsthm,mathtools}
\usepackage{microtype}
\usepackage[dvipsnames]{xcolor}
\usepackage{graphicx}
\usepackage{booktabs,longtable,array,tabularx}
\usepackage{enumitem}
\usepackage{fancyhdr}
\usepackage{titlesec}
\usepackage[most]{tcolorbox}
\usepackage{listings}
\usepackage{tikz}
\usetikzlibrary{positioning}
\usepackage{hyperref}
\hypersetup{colorlinks=true,linkcolor=MidnightBlue,urlcolor=RoyalBlue,pdfauthor={Independent Mathematics and Engineering Study Suite},pdftitle={PDE and Boundary Value Problems Guide}}
\definecolor{ink}{HTML}{162639}
\definecolor{navy}{HTML}{173B61}
\definecolor{blue}{HTML}{2874A6}
\definecolor{teal}{HTML}{007C83}
\definecolor{amber}{HTML}{8B5E16}
\definecolor{pale}{HTML}{EDF4F8}
\definecolor{paleteal}{HTML}{ECF8F5}
\definecolor{paleamber}{HTML}{FFF6E8}
\titleformat{\chapter}[display]{\normalfont\huge\bfseries\color{navy}}{\filleft\fontsize{43}{46}\selectfont\thechapter}{0.4ex}{\titlerule\vspace{1ex}}[\vspace{1ex}\titlerule]
\titleformat{\section}{\normalfont\Large\bfseries\color{navy}}{\thesection}{0.8em}{}
\titleformat{\subsection}{\normalfont\large\bfseries\color{teal}}{\thesubsection}{0.75em}{}
\pagestyle{fancy}\fancyhf{}\fancyhead[L]{\small\color{navy}Differential Equations Engineering Suite}\fancyhead[R]{\small\color{navy}Volume 6 of 8}\fancyfoot[C]{\thepage}\renewcommand{\headrulewidth}{0.3pt}
\setlist[itemize]{leftmargin=*,itemsep=3pt,topsep=5pt}\setlist[enumerate]{leftmargin=*,itemsep=4pt,topsep=5pt}
\setlength{\parskip}{5pt}\setlength{\parindent}{0pt}\setlength{\emergencystretch}{2em}
\newtheorem{theorem}{Theorem}[chapter]
\newtheorem{proposition}[theorem]{Proposition}
\newtheorem{lemma}[theorem]{Lemma}
\theoremstyle{definition}\newtheorem{definition}[theorem]{Definition}
\newtheorem{example}[theorem]{Worked example}
\theoremstyle{remark}\newtheorem{remark}[theorem]{Remark}
\newtcolorbox{bookground}[1][]{enhanced,breakable,colback=pale,colframe=blue,boxrule=0.6pt,title={TEXTBOOK ALIGNMENT},fonttitle=\sffamily\bfseries\small,#1}
\newtcolorbox{modern}[1][]{enhanced,breakable,colback=paleteal,colframe=teal,boxrule=0.6pt,title={MODERN COMPUTATIONAL EXTENSION},fonttitle=\sffamily\bfseries\small,#1}
\newtcolorbox{caution}[1][]{enhanced,breakable,colback=paleamber,colframe=amber,boxrule=0.6pt,title={ASSUMPTION / VALIDATION},fonttitle=\sffamily\bfseries\small,#1}
\newcommand{\dd}{\,\mathrm{d}}
\newcommand{\R}{\mathbb{R}}
\newcommand{\lap}{\mathcal{L}}
\newcommand{\norm}[1]{\left\lVert#1\right\rVert}
\newcommand{\abs}[1]{\left\lvert#1\right\rvert}
\lstdefinestyle{engpy}{language=Python,basicstyle=\ttfamily\footnotesize,keywordstyle=\color{RoyalBlue}\bfseries,commentstyle=\color{gray},stringstyle=\color{teal},backgroundcolor=\color{pale},frame=single,rulecolor=\color{blue!30},breaklines=true,columns=fullflexible,keepspaces=true,showstringspaces=false,tabsize=4,upquote=true}
\lstset{style=engpy}
\begin{document}
\begin{titlepage}
\pagecolor{navy}\color{white}
\vspace*{0.75in}
{\sffamily\bfseries\large MATHEMATICS \& ENGINEERING DOCUMENTATION SUITE\\[12pt]}
{\color{cyan!40}\rule{\textwidth}{2pt}}\\[1.0in]
{\sffamily\bfseries\fontsize{29}{36}\selectfont PDE and Boundary Value Problems Guide\par}
\vspace{0.45in}
{\Large Volume 6 of 8}\par\vspace{0.55in}
{\color{cyan!30}\rule{0.50\textwidth}{1pt}}\par\vspace{0.3in}
{\large Independent technical companion aligned with the\par}
{\large 13-chapter organization of Nagle, Saff \& Snider (7th ed.)\par}
\vfill
{\large Original derivations, worked examples, verification and engineering practice\par}\vspace{0.3in}
{\small Prepared October 2026 \quad\textbullet\quad pdf\LaTeX{} compatible\par}
\end{titlepage}
\nopagecolor\color{ink}\pagenumbering{roman}
\chapter*{Reader's guide and source boundaries}\addcontentsline{toc}{chapter}{Reader's guide and source boundaries}
This independent study resource is organized around \emph{Fundamentals of Differential Equations and Boundary Value Problems}, Seventh Edition, by R.~Kent Nagle, Edward~B. Saff, and Arthur~David Snider (Pearson, 2018; ISBN 978-0-321-97710-6). The book's chapter and section numbers are given as \emph{navigation references}, not as claims of reproduced text or proofs.

\begin{bookground}
\textbf{Textbook-grounded scope:} topics, terminology, and chapter/section organization corresponding to the cited textbook. All explanations, derivations, exercise prompts, and diagrams in this volume are newly authored and may emphasize different examples. Check the textbook itself for its exact hypotheses, proofs, exercises, or nomenclature in a particular section.
\end{bookground}
\begin{modern}
\textbf{Distinct extension layer:} Python/NumPy/SciPy workflows, engineering QA and reproducibility practices, automated convergence studies, diagnostic plots, and implementation guidance. These are \emph{not} represented as content supplied by the textbook. API calls are examples and should be verified against the installed package version.
\end{modern}
\begin{caution}
\textbf{Mathematical qualification:} formulas require their stated regularity, domain and boundary/initial conditions. Numerical agreement does not replace a proof; computational outputs are not certified error bounds unless explicitly established.
\end{caution}
\tableofcontents\clearpage\pagenumbering{arabic}
\chapter{From physical conservation to PDEs}
\begin{bookground}Correspondence: Chapter 10 (Partial Differential Equations), Sections 10.1--10.7; Chapter 11 (Eigenvalue Problems and Sturm--Liouville Equations), Sections 11.1--11.8.\end{bookground}
\section{Parabolic, hyperbolic, and elliptic prototypes}
The \emph{heat equation} $u_t=\kappa u_{xx}$ models diffusion for $\kappa>0$; the \emph{wave equation} $u_{tt}=c^2u_{xx}$ models idealized undamped propagation; \emph{Laplace's equation} $\Delta u=0$ models steady potential fields. These classifications describe principal parts under appropriate assumptions and are not interchangeable with boundary/initial data.
\section{Initial and boundary conditions}
A heat equation on $0<x<L$ needs one initial temperature profile and appropriate spatial boundary conditions. A wave equation needs initial displacement and initial velocity, plus spatial boundary data. An elliptic BVP generally uses spatial boundary conditions without an initial time. For each model specify domain, coefficient values and regularity of forcing and data.
\section{Conservation interpretation}
For heat transport with conductivity $k$, density $\rho$ and heat capacity $c_p$, Fourier's law $q=-k\nabla T$ and energy balance yield $\rho c_p T_t=\nabla\cdot(k\nabla T)+Q$. The textbook's constant-coefficient one-dimensional heat equation is recovered when material parameters are uniform and $Q=0$, with diffusivity $\kappa=k/(\rho c_p)$.
\chapter{Separation of variables and Fourier series}
\begin{bookground}Correspondence: Chapter 10, Sections 10.2--10.4.\end{bookground}
\section{The separated ansatz}
For $u_t=\kappa u_{xx}$ with homogeneous Dirichlet conditions, seek $u=X(x)T(t)$, obtaining $T'/(\kappa T)=X''/X=-\lambda$. The spatial eigenproblem is $X''+\lambda X=0$, $X(0)=X(L)=0$. Its nontrivial modes have
\[\lambda_n=(n\pi/L)^2,\qquad X_n(x)=\sin(n\pi x/L),\quad n\ge1.\]
The corresponding time factors are $T_n(t)=e^{-\kappa\lambda_n t}$.
\section{Fourier sine and cosine expansions}
An initial profile $f$ with appropriate regularity has sine coefficients
\[b_n=\frac2L\int_0^L f(x)\sin\frac{n\pi x}{L}\dd x,\]
leading to $u(x,t)=\sum_{n\ge1}b_n e^{-\kappa(n\pi/L)^2t}\sin(n\pi x/L)$. Neumann conditions yield cosine modes including a constant $n=0$ mode. Convergence at endpoints and discontinuities follows the relevant Fourier theory; a truncated series may exhibit Gibbs oscillations near jumps.
\begin{example}[A single thermal mode]
With $u(x,0)=3\sin(2\pi x/L)$ and fixed-zero endpoints, all coefficients vanish except $b_2=3$, so
\[u(x,t)=3e^{-4\kappa\pi^2t/L^2}\sin(2\pi x/L).\]
The solution can be checked by differentiating and using $u_{xx}=-(2\pi/L)^2u$.
\end{example}
\chapter{Heat, wave and Laplace equations}
\begin{bookground}Correspondence: Chapter 10, Sections 10.5 (heat), 10.6 (wave), and 10.7 (Laplace).\end{bookground}
\section{Heat-equation maximum principle}
Under classical regularity and standard Dirichlet conditions for the homogeneous heat equation, interior maxima over a forward space-time cylinder are controlled by initial and boundary data. Numerical temperatures violating physical bounds may signal unstable discretization, incorrectly handled boundary data, or an inappropriate physical source model.
\section{Wave equation modal solution}
For fixed endpoints and initial displacement $f$, velocity $g$,
\[u(x,t)=\sum_{n=1}^{\infty}\left(A_n\cos\frac{n\pi ct}{L}+B_n\sin\frac{n\pi ct}{L}\right)\sin\frac{n\pi x}{L},\]
where $A_n$ are sine coefficients of $f$ and $B_n=(L/(n\pi c))$ times sine coefficients of $g$. Solutions depend on the strength and regularity of both initial fields. Ideal strings without damping do not decay in energy.
\section{Laplace equation and steady states}
For $u_{xx}+u_{yy}=0$ in a rectangle, separated solutions use trigonometric and hyperbolic factors chosen to satisfy particular edge conditions. In polar geometry, angular Fourier modes combine with radial powers; behavior at the origin or infinity constrains admissible solutions. Boundary value geometry dictates the basis.
\chapter{Sturm--Liouville eigenvalue problems}
\begin{bookground}Correspondence: Chapter 11, Sections 11.1--11.3 and 11.7--11.8. The book's canonical form is $(py')'+qy+\lambda r y=0$ with positive weight $r$ for regular problems.\end{bookground}
\section{Regular self-adjoint form}
Write
\[-(p y')'-q y=\lambda r y\quad\text{on }(a,b),\]
with suitable regular coefficients $p>0$ and weight $r>0$, and separated homogeneous endpoint conditions making the boundary form vanish. A regular self-adjoint problem has real eigenvalues and orthogonal eigenfunctions for distinct eigenvalues. Under the standard regular separated setting, eigenvalues form an increasing unbounded sequence and eigenfunctions provide a powerful expansion system.
\section{Orthogonality derivation}
Suppose $Ly_m=\lambda_m r y_m$ and $Ly_n=\lambda_n r y_n$ with $L=-(p\partial_x)'-q$. Multiply the first equation by $y_n$, the second by $y_m$, subtract and integrate:
\[(\lambda_m-\lambda_n)\int_a^b r y_my_n\dd x=\left[p(y_m y_n'-y_n y_m')\right]_a^b.\]
For self-adjoint separated boundary conditions the boundary form vanishes, hence $\lambda_m\ne\lambda_n$ implies weighted orthogonality. This proof relies on the correct domain and endpoint conditions.
\section{Rayleigh quotient}
For a self-adjoint boundary problem with boundary term eliminated by appropriate conditions,
\[\lambda=\frac{\int_a^b[p(y')^2-q y^2]\dd x}{\int_a^b r y^2\dd x}\]
for a nontrivial eigenfunction. A trial function satisfying admissible boundary constraints can supply variational estimates under the usual minimax hypotheses. Do not use this formula unchanged when boundary terms remain or coefficients/conditions differ.
\section{Singular problems and special functions}
When $p$ or $r$ degenerates at an endpoint or the interval is unbounded, regular spectral conclusions may not transfer directly. The textbook discusses singular Sturm--Liouville problems and equations associated with Bessel, Legendre and Hermite functions; the relevant endpoint behavior and admissibility criteria must be investigated separately.
\chapter{Nonhomogeneous BVPs and Fredholm alternative}
\begin{bookground}Correspondence: Chapter 11, Section 11.4 (Fredholm Alternative), 11.5 (Eigenfunction Expansion), 11.6 (Green's Functions).\end{bookground}
\section{Resonance and compatibility}
For $-y''=f$ on $0<x<\pi$ with zero Dirichlet endpoints, the operator is invertible and a unique solution exists for suitable $f$. For $y''+y=f$ on the same interval with the same endpoints, the homogeneous equation has the eigenfunction $\sin x$ in its kernel, and compatibility requires $\int_0^\pi f(x)\sin x\dd x=0$. When compatible, homogeneous multiples of $\sin x$ can be added, so uniqueness fails without an additional condition.
\section{Eigenfunction expansion}
Suppose $Ly_n=\lambda_n r y_n$ with eigenfunctions normalized in weighted $L^2$ and suppose $Ly=f$ has a solvable expansion. Formally, $y=\sum_n c_n y_n$ with
\[c_n=\frac{\langle f/r,y_n\rangle_r}{\lambda_n}=\frac{\int_a^b f(x)y_n(x)\dd x}{\lambda_n}\]
for $\lambda_n\ne0$ under this operator convention. Zero eigenvalues require compatibility. Convergence and boundary regularity demand separate justification; an expansion is not automatically pointwise valid everywhere.
\section{Green's functions}
If a linear BVP admits an inverse, its solution may be represented as $y(x)=\int_a^bG(x,\xi)f(\xi)\dd\xi$. For $-y''=f$ with $y(0)=y(L)=0$, a suitable Green's function is
\[G(x,\xi)=\begin{cases}x(L-\xi)/L,&0\le x\le\xi,\\\xi(L-x)/L,&\xi\le x\le L.\end{cases}\]
It is continuous at $x=\xi$, satisfies both endpoint conditions, and its first derivative has jump $G_x(\xi^+,\xi)-G_x(\xi^-,\xi)=-1$, as required by $-G_{xx}=\delta(x-\xi)$.
\begin{example}[Green's function benchmark]
For $-y''=2$ and zero Dirichlet endpoints on $[0,1]$, the Green integral gives $y(x)=x(1-x)$. It satisfies $-y''=2$ and both endpoint constraints.
\end{example}
\chapter{Numerical treatment and validation}
\begin{modern}Finite differences, collocation with SciPy and sparse eigenvalue algorithms are modern implementation pathways corresponding to ideas in the textbook's Chapter 10 projects and Chapter 11 project sections. The specific APIs below are external extensions.\end{modern}
\section{Finite-difference eigenvalue problem}
Let $h=L/(N+1)$ and apply $-y''\approx(-y_{i-1}+2y_i-y_{i+1})/h^2$. The resulting symmetric positive-definite matrix approximates the Dirichlet Laplacian. Compare its lowest eigenvalues with $(n\pi/L)^2$, refine $N$, and normalize eigenvectors consistently before comparing. The eigenfunction sign is arbitrary.
\begin{lstlisting}
import numpy as np
from scipy.sparse import diags
from scipy.sparse.linalg import eigsh
L, N = 1.0, 100
h = L/(N+1)
A = diags([-np.ones(N-1), 2*np.ones(N), -np.ones(N-1)],
          [-1, 0, 1], format='csr') / h**2
lam, vec = eigsh(A, k=2, which='SM')
lam = np.sort(lam)
assert np.allclose(lam, (np.arange(1,3)*np.pi)**2, rtol=0.003)
\end{lstlisting}
\section{Nonhomogeneous boundary residuals}
When using numerical BVP solvers, record the maximum defect in differential equations (using suitably accurate differentiation or collocation residual evaluation) \emph{and} each endpoint residual. The latter are especially important at or near eigenvalues, where an apparently small discretization residual may coexist with very large solution sensitivity.
\section{Heat equation energy check}
For zero Dirichlet heat flow, $E(t)=\frac12\int_0^L u(x,t)^2\dd x$ satisfies $E'(t)=-\kappa\int_0^L u_x^2\dd x\le0$ under classical smoothness assumptions. This supplies an independent validation invariant for finite-difference and method-of-lines implementations.
\chapter{Engineering exercises and solution guidance}
\begin{enumerate}
\item Find eigenvalues/eigenfunctions of $-y''=\lambda y$ on $(0,L)$ with zero Dirichlet endpoints.
\item Solve $u_t=\kappa u_{xx}$ on $(0,L)$ with $u(x,0)=\sin(3\pi x/L)$ and zero endpoints.
\item Demonstrate a nonunique solution family for $y''+y=0$ with $y(0)=y(\pi)=0$.
\item Derive the Green kernel for $-y''$ with Dirichlet endpoints from piecewise linear solutions and a unit derivative jump.
\item Explain why Neumann heat flow preserves total thermal energy in the insulated, unforced, constant-property one-dimensional model.
\item Perform a mesh-refinement study for the first Dirichlet eigenvalue and report $|\lambda_{1,h}-\pi^2|$.
\end{enumerate}
\section{Answer checks}
Problems 1 and 2 give $\lambda_n=(n\pi/L)^2$, $y_n=\sin(n\pi x/L)$, and $u(x,t)=e^{-9\kappa\pi^2t/L^2}\sin(3\pi x/L)$. Problem 3 has $y=C\sin x$ for arbitrary $C$. Problem 4 yields the piecewise kernel stated above. Problem 5 follows by integrating $u_t=\kappa u_{xx}$ and applying $u_x(0,t)=u_x(L,t)=0$. For problem 6, the standard central-difference eigenvalue is $\lambda_{1,h}=4h^{-2}\sin^2(\pi/[2(N+1)])$ with $h=1/(N+1)$.

\appendix
\chapter{References and reproducibility}
\begin{enumerate}
\item R.~K. Nagle, E.~B. Saff, and A.~D. Snider, \emph{Fundamentals of Differential Equations and Boundary Value Problems}, 7th ed., Pearson, 2018, ISBN 978-0-321-97710-6. Primary source for the \emph{topic map and terminology}; cited locally by chapter and section.
\item SciPy project, \emph{solve\_ivp} reference, \url{https://docs.scipy.org/doc/scipy/reference/generated/scipy.integrate.solve_ivp.html} (modern extension).
\item SciPy project, \emph{solve\_bvp} reference, \url{https://docs.scipy.org/doc/scipy/reference/generated/scipy.integrate.solve_bvp.html} (modern extension).
\item SciPy project, \emph{eigsh} reference, \url{https://docs.scipy.org/doc/scipy/reference/generated/scipy.sparse.linalg.eigsh.html} (modern extension).
\item Python, NumPy, SciPy, and SymPy project documentation; computational checks require recording package versions and tolerances.
\end{enumerate}
\section*{Verification record template}
\begin{tabularx}{\textwidth}{>{\bfseries}p{0.30\textwidth}X}\toprule
Problem identifier & \rule{0pt}{15pt}\ \\ \midrule
ODE/PDE and domain & \\[8pt]
Initial/boundary conditions & \\[8pt]
Units and parameter values & \\[8pt]
Analytic residual & \\ [8pt]
Grid/step refinement & \\ [8pt]
Independent solver comparison & \\ [8pt]
Version and reproducibility notes & \\ [8pt]\bottomrule
\end{tabularx}
\end{document}
